% -----------------------------------------------------------
\section{Damnation, Damn}
% -----------------------------------------------------------
	\label{DAMNATION}
	
% % ----------------------
% \subsection{See also}
% % ----------------------

% ----------------------
\subsection{1828 American Dictionary of the English Language} % *1828
% ----------------------
	\label{DAMNATION-1828}
	
		% -----
		\subsubsection{Damnation}
		% ----- 

			\begin{compactitem}
				\item[1] Sentence or condemnation to everlasting punishment in the future state; or the state of eternal torments.
				\item[2] Condemnation.
			\end{compactitem}
			
		% -----
		\subsubsection{Damn}
		% ----- 

			\begin{compactitem}
				\item[1] To sentence to eternal torments in a future state; to punish in hell.
				\item[2] To condemn; to decide to be wrong or worthy of punishment; to censure; to reprobate.
				\item[3] To condemn; to explode; to decide to be bad, mean, or displeasing, be hissing or any mark of disapprobation; as, to \textit{damn} a play, or a mean author.
				\item[4] A word used in profaneness; a term of execration.
			\end{compactitem}

% % ----------------------
% \subsection{Oxford English Dictionary} % *OED
% % ----------------------
	% \label{DAMNATION-OED}

	% \begin{compactitem}
		% \item[]
	% \end{compactitem}

% ----------------------
\subsection{Scriptural Usage} % *SCRIPTURES
% ----------------------
	\label{DAMNATION-SCRIPTURES}

	% % -----
	% \subsubsection{Old Testament} % *OT
	% % -----
		% \label{DAMNATION-OT}
	
		% % --
		% \paragraph{KJV}
		% % --
			% \label{DAMNATION-OT-KJV}
	
			% \begin{tabular}{ll}
				% Total occurrences in the OT & 
			% \end{tabular}
			
			% \begin{compactdesc}
				% \item[]
			% \end{compactdesc}
			
		% % --
		% \paragraph{Hebrew Bible}
		% % --
			% \label{DAMNATION-HEBREW}
		
			% %
			% \subparagraph{\texorpdfstring{\he{sample word}}{}}
			% %
				% \label{PAGA} % whatever is in step bible without periods
			
				% \begin{compactdesc}
					% \item[HALOT , PDF ] \textbf{}
						% \begin{compactitem}
							% \item[1]
						% \end{compactitem}
					% \item[BDB , PDF ] \textbf{}
						% \begin{compactitem}
							% \item[1]
						% \end{compactitem}
					% \item[DCH PDF ] \textbf{}
						% \begin{compactitem}
							% \item[1]
						% \end{compactitem}
				% \end{compactdesc}
				
				% \begin{compactdesc}
					% \item[Some OT uses] \textbf{}
						% \begin{compactdesc}
							% \item[]
						% \end{compactdesc}
				% \end{compactdesc}
			
		% % --
		% \paragraph{Septuagint}
		% % --
			% \label{DAMNATION-LXX}
		
			% %
			% \subparagraph{\texorpdfstring{\gr{sample word}}{}}
			% %
		
	% -----
	\subsubsection{New Testament} % *NT
	% -----
		\label{DAMNATION-NT}
	
		% --
		\paragraph{KJV}
		% --
			\label{DAMNATION-NT-KJV}
			
			See \hyperref[DAMNATION-new-testament]{below} for some notes on these uses.
	
			\begin{tabular}{ll}
				Total occurrences in the NT & 15
			\end{tabular}
			
			\begin{compactdesc}
				\item[Matthew 23:14] Woe unto you, scribes and Pharisees, hypocrites!  for ye devour widows' houses, and for a pretence make long prayer: therefore ye shall receive the greater damnation.
				\item[Matthew 23:33] Ye serpents, ye generation of vipers, how can ye escape the damnation of hell?
				\item[Mark 3:29] But he that shall blaspheme against the Holy Ghost hath never forgiveness, but is in danger of eternal damnation:
				\item[Mark 12:40] Which devour widows' houses, and for a pretence make long prayers: these shall receive greater damnation.
				\item[Mark 16:16] He that believeth and is baptized shall be saved; but he that believeth not shall be damned.
				\item[Luke 20:47] Which devour widows' houses, and for a shew make long prayers: the same shall receive greater damnation.
				\item[John 5:29] And shall come forth; they that have done good, unto the resurrection of life; and they that have done evil, unto the resurrection of damnation.
				\item[Romans 3:8] And not rather, (as we be slanderously reported, and as some affirm that we say,) Let us do evil, that good may come?  whose damnation is just.
				\item[Romans 13:2] Whosoever therefore resisteth the power, resisteth the ordinance of God: and they that resist shall receive to themselves damnation.
				\item[Romans 14:23] And he that doubteth is damned if he eat, because he eateth not of faith: for whatsoever is not of faith is sin.
				\item[1 Corinthians 11:29] For he that eateth and drinketh unworthily, eateth and drinketh damnation to himself, not discerning the Lord's body.
				\item[2 Thessalonians 2:12] That they all might be damned who believed not the truth, but had pleasure in unrighteousness.
				\item[1 Timothy 5:12] Having damnation, because they have cast off their first faith.
				\item[2 Peter 2:1] BUT there were false prophets also among the people, even as there shall be false teachers among you, who privily shall bring in damnable heresies, even denying the Lord that bought them, and bring upon themselves swift destruction.
				\item[2 Peter 2:3] And through covetousness shall they with feigned words make merchandise of you: whose judgment now of a long time lingereth not, and their damnation slumbereth not.
			\end{compactdesc}
			
		% --
		\paragraph{Greek New Testament} % *GREEK
		% --
			\label{DAMNATION-GREEK}
		
			%
			\subparagraph{\texorpdfstring{\gr{κρίσις}}{}}
			%
				\label{KRISIS}
			
				\begin{compactdesc}
					\item[BDAG 504a] \textbf{}
						\begin{compactitem}
							\item[1] legal process of judgment, \textit{judging, judgment}
							\item[1a] of the activity of God or the Messiah as judge, especially on the Last Day
							\item[1a.$\alpha$] God's righteous \textit{judgment}
							\item[1a.$\beta$] The word often means \textit{judgment} that goes against a person, \textit{condemnation}, and the \textit{sentence} that follows
							\item[1b] of the \textit{judgment} of one person upon or against another, in the nature of an \textit{evaluation}
							\item[2] a board of judges, \textit{court}
							\item[3] administration of what is right and fair, \textit{right} in the sense of \textit{justice/righteousness}
						\end{compactitem}
					\item[LSJ 997a, PDF 1068] \textbf{}
						\begin{compactitem}
							\item[I.1] \textit{separating, distinguishing}
							\item[I.2] \textit{decision, judgement}
							\item[I.3] \textit{choice, election}
							\item[I.4] \textit{interpretation} of dreams or portents
							\item[II.1a] \textit{judgement} of a court
							\item[II.1b] \textit{result of a trial, condemnation}
							\item[II.2] \textit{trial} of skill or strength
							\item[II.3] \textit{dispute}
							\item[III.1] \textit{event, issue}
							\item[III.2] \textit{turning point} of a disease, \textit{sudden change} for better or worse
							\item[IV] middle of the spinal column
						\end{compactitem}
				\end{compactdesc}
				
				\begin{compactdesc}
					\item[Some NT uses] \textbf{}
						\begin{compactdesc}
							\item[Matthew 23:33] \gr{ὄφεις γεννήματα ἐχιδνῶν, πῶς φύγητε ἀπὸ τῆς κρίσεως τῆς γεέννης;} \textit{(This is going to be BDAG \#1a.$\alpha$, I'm suppose; the negative connotation is the key I guess?)}
							\item[Mark 3:29] \gr{ὃς δ᾽ ἂν βλασφημήσῃ εἰς τὸ πνεῦμα τὸ ἅγιον, οὐκ ἔχει ἄφεσιν εἰς τὸν αἰῶνα, ἀλλὰ ἔνοχός ἐστιν αἰωνίου ἁμαρτήματος.} \textit{(There is an alternate text where you get \gr{κρίσεως} instead of \gr{ἁμαρτήματος}.)}
							\item[2 Thessalonians 1:5] \gr{νδειγμα τῆς δικαίας κρίσεως τοῦ θεοῦ, εἰς τὸ καταξιωθῆναι ὑμᾶς τῆς βασιλείας τοῦ θεοῦ, ὑπὲρ ἧς καὶ πάσχετε,}
							\item[John 5:29] \gr{καὶ ἐκπορεύσονται οἱ τὰ ἀγαθὰ ποιήσαντες εἰς ἀνάστασιν ζωῆς, οἱ δὲ τὰ φαῦλα πράξαντες εἰς ἀνάστασιν κρίσεως.}
						\end{compactdesc}
				\end{compactdesc}
		
			%
			\subparagraph{\texorpdfstring{\gr{κρίμα}}{}}
			%
				\label{KRIMA}
			
				\begin{compactdesc}
					\item[BDAG 502a] \textbf{}
						\begin{compactitem}
							\item[1] legal action taken against someone, \textit{dispute, lawsuit}
							\item[2] content of a deliberative process, \textit{decision, decree}
							\item[3] action or function of a judge, \textit{judging, judgment}
							\item[4] legal decision rendered by a judge, \textit{judicial verdict}
							\item[4a] generally
							\item[4b] mostly in an unfavorable sense, of the \textit{condemnatory verdict} and sometimes the subsequent \textit{punishment} itself
							\item[5] proper recognition of someone's rights, \textit{justice}
							\item[6] in John, \gr{κρίμα} shows the same two-sidedness as the other members of the \gr{κρίνω} family (``judgment'' and ``separation''), and means the judicial decision which consists in the separation of those wh	o are willing to believe from those who are unwilling to do so. \textit{(Note that the separation being a ``judicial decision'' suggests that it is based in} laws\textit{, which is what you would expect.)}
						\end{compactitem}
					\item[LSJ 995b, PDF 1066] \textbf{}
						\begin{compactitem}
							\item[I.1] \textit{decision, judgement}
							\item[I.2] \textit{decree, resolution}
							\item[I.3] \textit{legal decision}
							\item[II.1] \textit{matter for judgement, question}
							\item[II.2] \textit{law-suit}
							\item[III.1] = \gr{κρίσις}, \textit{judging, judgement}
						\end{compactitem}
				\end{compactdesc}
				
				\begin{compactdesc}
					\item[Some NT uses] \textbf{}
						\begin{compactdesc}
							\item[Matthew 23:14] \gr{Οὐαὶ ὑμῖν, γραμματεῖς καὶ Φαρισαῖοι ὑποκριταί, ὅτι κατεσθίετε τὰς οἰκίας τῶν χηρῶν καὶ προφάσει μακρὰ προσευχόμενοι· διὰ τοῦτο λήψεσθε περισσότερον κρίμα} \textit{(This verse is removed in the critical edition.)}
							\item[Mark 12:40] \gr{οἱ κατεσθίοντες τὰς οἰκίας τῶν χηρῶν καὶ προφάσει μακρὰ προσευχόμενοι· οὗτοι λήμψονται περισσότερον κρίμα.}
							\item[Luke 20:47] \gr{οἳ κατεσθίουσιν τὰς οἰκίας τῶν χηρῶν καὶ προφάσει μακρὰ προσεύχονται· οὗτοι λήμψονται περισσότερον κρίμα.}
							\item[Romans 3:8] \gr{καὶ μὴ καθὼς βλασφημούμεθα καὶ καθώς φασίν τινες ἡμᾶς λέγειν ὅτι Ποιήσωμεν τὰ κακὰ ἵνα ἔλθῃ τὰ ἀγαθά; ὧν τὸ κρίμα ἔνδικόν ἐστιν.}
							\item[Romans 13:2] \gr{ὥστε ὁ ἀντιτασσόμενος τῇ ἐξουσίᾳ τῇ τοῦ θεοῦ διαταγῇ ἀνθέστηκεν, οἱ δὲ ἀνθεστηκότες ἑαυτοῖς κρίμα λήμψονται.}
							\item[1 Corinthians 11:29] \gr{ὁ γὰρ ἐσθίων καὶ πίνων κρίμα ἑαυτῷ ἐσθίει καὶ πίνει μὴ διακρίνων τὸ σῶμα.}
							\item[1 Timothy 5:12] \gr{ἔχουσαι κρίμα ὅτι τὴν πρώτην πίστιν ἠθέτησαν·}
						\end{compactdesc}
				\end{compactdesc}
				
			%
			\subparagraph{\texorpdfstring{\gr{κατακρίνω}}{}}
			%
				\label{KATAKRINO}
			
				\begin{compactdesc}
					\item[BDAG 460a] \textbf{}
						\begin{compactitem}
							\item pronounce a sentence after determination of guilt, \textit{pronounce a sentence on} someone
						\end{compactitem}
					\item[LSJ 896a, PDF 967] \textbf{}
						\begin{compactitem}
							\item[I.1] \textit{give as sentence against}
							\item[I.2] \textit{condemn}
							\item[II] in passive, simply \textit{to be judged, deemed}
						\end{compactitem}
				\end{compactdesc}
				
				\begin{compactdesc}
					\item[Some NT uses] \textbf{}
						\begin{compactdesc}
							\item[Mark 16:16] \gr{ὁ πιστεύσας καὶ βαπτισθεὶς σωθήσεται, ὁ δὲ ἀπιστήσας κατακριθήσεται.}
							\item[Romans 14:23] \gr{ὁ δὲ διακρινόμενος ἐὰν φάγῃ κατακέκριται, ὅτι οὐκ ἐκ πίστεως· πᾶν δὲ ὃ οὐκ ἐκ πίστεως ἁμαρτία ἐστίν.}
						\end{compactdesc}
				\end{compactdesc}
		
	% % -----
	% \subsubsection{Book of Mormon} % *BOM
	% % -----
		% \label{DAMNATION-BOM}
	
		% \begin{tabular}{ll}
			% Total occurrences in the BOM & 
		% \end{tabular}
		
		% \begin{compactdesc}
			% \item[]
		% \end{compactdesc}
		
	% % -----
	% \subsubsection{Doctrine and Covenants} % *DC
	% % -----
		% \label{DAMNATION-DC}
	
		% \begin{tabular}{ll}
			% Total occurrences in the DC & 
		% \end{tabular}
		
		% \begin{compactdesc}
			% \item[]
		% \end{compactdesc}
		
	% % -----
	% \subsubsection{Pearl of Great Price} % *PGP
	% % -----
		% \label{DAMNATION-PGP}
	
		% \begin{tabular}{ll}
			% Total occurrences in the PGP & 
		% \end{tabular}
		
		% \begin{compactdesc}
			% \item[]
		% \end{compactdesc}
		
% % ----------------------
% \subsection{Key Phrases} % *KEYPHRASES *PHRASES
% % ----------------------
	% \label{DAMNATION-PHRASES}

	% \begin{compactdesc}
		% \item[]
	% \end{compactdesc}
	
% % ----------------------
% \subsection{Bible Dictionaries} % *DICTIONARIES
% % ----------------------
	% \label{DAMNATION-BD}

% % ----------------------
% \subsection{Restoration Usage} % *RESTORATION
% % ----------------------
	% \label{DAMNATION-RESTORATION}

	% % -----
	% \subsubsection{Joseph Smith} % *JS
	% % -----
		% \label{DAMNATION-JS}
	
		% \begin{compactdesc}
			% \item[DATE/SOURCE]
		% \end{compactdesc}
	
	% % -----
	% \subsubsection{Brigham Young} % *BY
	% % -----
		% \label{DAMNATION-BY}
	
		% \begin{compactdesc}
			% \item[DATE/SOURCE]
		% \end{compactdesc}
	
% ----------------------
\subsection{Commentary and Studies} % *COMMENTARY *STUDIES
% ----------------------
	\label{DAMNATION-COMMENTARY}

	% -----
	\subsubsection{Key Points}
	% -----
		\label{DAMNATION-KEYS}
	
		\begin{compactitem}
			\item There is little to no specific evidence in the NT for supporting the traditional LDS understanding of ``damnation'' as ``cessation of progress,'' at least when you look at the uses of ``damnation'' in the NT and their Greek basis. You have to support it some other way.
			\item The Greek words underlying ``damnation'' in the NT don't really carry with them any idea of eternal punishment; see \hyperref[DAMNATION-new-testament]{below}.
		\end{compactitem}
		
	% -----
	\subsubsection{Damnation in the New Testament}
	% -----
		\label{DAMNATION-new-testament}
		
		It's interesting to me that none of the words underlying ``damnation'' in the NT convey an idea of eternal punishment. They may---but definitely not in every case---have a sense that puts emphasis on a punishment, but they don't specify what that punishment is nor its duration. You would have to gather something about the nature of the punishment from context or from some other source entirely.
		
		Note that this is very different from the English word ``damnation,'' which as you can see from the \hyperref[DAMNATION-1828]{1828 definitions} above, clearly does connote the idea of an eternal punishment. What this means is that interpretations of KJV verses involving ``damnation'' will almost inevitably be too harsh, as we read the force of the English word into the passage, even though the Greek word does not suggest the same force.
		
		Of the two Greek words that I studied above, \hyperref[KRIMA]{\gr{κρίμα}} in meaning \#4b is the one closest to the English word ``damnation,'' but even that isn't really the same thing. The other word, \hyperref[KRISIS]{\gr{κρίσις}}, is a bit farther away, though meaning 1.a.$\beta$ is in the ballpark.
		
		This is important especially for a verse like John 5:29, which in the KJV reads as follows:
		
			\begin{quote}
				And shall come forth; they that have done good, unto the resurrection of life; and they that have done evil, unto the resurrection of damnation.
			\end{quote}
			
		Virtually all modern translations replace the word ``damnation'' at the end with something less harsh, like ``condemnation'' or even ``judgment.'' The Greek word is \gr{κρίσις} so it's the weaker of the two main Greek words, at least in relation to the idea of punishment.
		
		John 5:29 is a very important verse from the perspective of the restoration, because it's the one that JS and Rigdon were translating in their NT revision when they received the revelation in \hyperref[DC76]{DC 76}. There is no hint of ``damnation'' anywhere in that revelation, as it's almost all light and hope for everyone who is resurrected. The word ``damnation'' isn't even used in the discussion of sons of perdition early in the section.
		
	% -----
	\subsubsection{Damnation as cessation of progress}
	% -----
		\label{DAMNATION-progress}
		
		\begin{compactitem}
			\item \hyperref[DC58-29]{DC 58:29}
		\end{compactitem}